% SOURCE AUTHORITY:
% expose_VIII_body.tex lines 924--1021; scan indices 254--257; source folios 243--246.
% Direct scan comparison: physical PDF pages 255--258 at 1100 dpi.
% Transparent source dispositions:
% - the proof heading's printed 2.5.11 is restored as 2.3.11;
% - typewritten letter-o indices and cochain degrees are normalized to zero;
% - the arrow-key slip inside the defining formula (2.3.14) is restored as
%   the minus sign already printed in (2.3.13) and required by the cocycle.
% Hard stop before part c), source line 1022.

\medskip
\noindent\underline{Proof of 2.3.11.}\footnote{For a more elegant
proof, see 2.3.19 below.}
Let $\xi$ be an element of $\mathrm{Biext}^1(P,Q;G)$, let
\[
 \varphi_\xi:{}_nP\otimes{}_nQ\longrightarrow G
\]
be the associated pairing, and let
\[
 u_\xi:{}_nP\longrightarrow{}_n\uExt^1(nQ,G)
\]
be the image of $\xi$ under (2.3.5).  We must prove that $-u_\xi$ is
obtained from $\varphi_\xi$ by the arrow (2.3.9), that is, that it is
the composite
\begin{equation}\tag{*}
 {}_nP
 \xrightarrow{\ \varphi_\xi\ }
 \uHom({}_nQ,G)
 \xrightarrow{\ (2.3.8)\ }
 {}_n\uExt^1(nQ,G).
\end{equation}
To prove this, we may take a section $x$ of $\,{}_nP$ and prove that
its image under $-u_\xi$ equals its image under the composite (*).
If desired, this would allow us to reduce to the case
$P=(\mathbb Z/n\mathbb Z)_T$, but we shall not need that reduction.

\medskip
\noindent a) Choose flat resolutions $L_\bullet$ and $M_\bullet$,
\[
 0\longrightarrow L_1\longrightarrow L_0\longrightarrow0,
 \qquad
 0\longrightarrow M_1\longrightarrow M_0\longrightarrow0,
\]
of $P$ and $Q$, respectively.  A flat resolution of $P$, for example,
is obtained by writing $P$ as a quotient of a flat
$\mathbb Z_T$-module $L_0$ and taking $L_1$ to be the kernel of
$L_0\longrightarrow P$; this kernel is flat because it is
torsion-free.  We obtain the complex $L_\bullet\otimes M_\bullet$
\[
 0\longrightarrow L_1\otimes M_1
 \xrightarrow{d_2}
 L_1\otimes M_0+L_0\otimes M_1
 \xrightarrow{d_1}
 L_0\otimes M_0
 \longrightarrow0,
\]
whose sheaf $\underline H_1$ will serve as $\uTor_1(P,Q)$.

Let $L'_0$ (respectively, $M'_0$) be the subsheaf of $L_0$
(respectively, $M_0$) which is the inverse image of $\,{}_nP$
(respectively, $\,{}_nQ$).  Thus
\[
 0\longrightarrow L_1\longrightarrow L'_0\longrightarrow0
 \qquad
 \left(\text{respectively, }
 0\longrightarrow M_1\longrightarrow M'_0\longrightarrow0\right)
\]
is a flat resolution of $\,{}_nP$ (respectively, of $\,{}_nQ$).  The
homomorphism
\[
 {}_nP\otimes{}_nQ
 \longrightarrow
 \uTor_1(P,Q)=\underline H_1(L_\bullet\otimes M_\bullet)
\]
may then be described, by passage to the quotient, from the
homomorphism
\[
 L'_0\otimes M'_0
 \longrightarrow
 \underline H_1(L_\bullet\otimes M_\bullet)
\]
given by
\[
 x_0\otimes y_0
 \longmapsto
 \text{the class of the $1$-cycle }
 [nx_0]\otimes y_0-x_0\otimes[ny_0].
\]
Here the brackets indicate that the element in question is regarded
as a section of the subsheaf $L_1$ of $L_0$ (respectively, $M_1$ of
$M_0$).  This description follows from the calculation made in the
proof of 2.1.10.

\medskip
\noindent b) Take an injective resolution
\[
 C^\bullet=\bigl(0\longrightarrow C^0
 \xrightarrow{d^0}C^1\longrightarrow\cdots\bigr)
\]
of $G$, and identify
\[
 \mathrm{Biext}^1(P,Q;G)
 \simeq
 \operatorname{Ext}^1\!\left(
   P\mathbin{\overset{\mathbb L}{\otimes}}Q,G
 \right)
 =H^1\!\left(
   \operatorname{Hom}^{\boldsymbol\cdot}
   (L_\bullet\otimes M_\bullet,C^{\boldsymbol\cdot})
 \right).
\]
The element $\xi$ may therefore be represented by a pair of
homomorphisms $f_1,f_0$, giving the anticommutative diagram
\begin{equation}\tag{2.3.12}
\begin{tikzcd}[column sep=huge,row sep=huge]
 L_1\otimes M_0+L_0\otimes M_1
   \arrow[r,"d_1"]
   \arrow[d,"f_1"']
   \arrow[rd,phantom,"-" description]&
 L_0\otimes M_0
   \arrow[d,"f_0"']\\
 C^0\arrow[r,"d^0"]&C^1.
\end{tikzcd}
\end{equation}
By definition, the pairing
\[
 \varphi_\xi:{}_nP\otimes{}_nQ\longrightarrow G
\]
is obtained, by passage to the quotient, from the homomorphism
\[
 \varphi:L'_0\otimes M'_0\longrightarrow G
\]
given by
\begin{equation}\tag{2.3.13}
 \varphi(x_0\otimes y_0)
 =f_1\bigl([nx_0]\otimes y_0-x_0\otimes[ny_0]\bigr)
 \in G=\ker(C^0\longrightarrow C^1).
\end{equation}
Since assertion 2.3.11 is local on $T$, we may suppose that the given
section $x$ of $\,{}_nP$ lifts to a section $x_0$ of $L'_0$.
Regarding $y_0$ as variable, the preceding formula then defines a
homomorphism
\begin{equation}\tag{2.3.14}
 \varphi':M'_0\longrightarrow G,
 \qquad
 y_0\longmapsto
 f_1\bigl([nx_0]\otimes y_0-x_0\otimes[ny_0]\bigr),
\end{equation}
which passes to the quotient to give
\[
 {}_nQ\longrightarrow G,
 \qquad
 y\longmapsto\varphi_\xi(x,y),
\]
the homomorphism associated with $\varphi_\xi$ and $x$.
